%% ============================================================
%% Teil I: C++ kompakt – was rclcpp wirklich voraussetzt
%% Band 3 – Kapitel 1 bis 4
%% ============================================================

\part{C++ kompakt}

% ─────────────────────────────────────────────────────────────
\chapter{Wir wollen C++ kompilieren – CMake und das ROS2-Build-System}
% ─────────────────────────────────────────────────────────────

\begin{lernziele}
Du verstehst den Unterschied zwischen kompilierter und
interpretierter Sprache, kannst ein C++-Programm mit g++
und mit CMake bauen und weißt, was ament\_cmake für ROS2 tut.
\end{lernziele}

\section{Situation}

In Python schreibst du \texttt{python3 sensor\_bridge.py} -- fertig.
In C++ schreibst du \texttt{motor\_node.cpp},
dann \texttt{g++ motor\_node.cpp -o motor\_node},
dann \texttt{./motor\_node}.
Und für ROS2 kommt CMake noch dazu.
Warum so viele Schritte?

\begin{aufgabeBox}
Kompiliere ein minimales C++-Programm mit g++ direkt,
dann mit CMake, und verstehe, was \texttt{CMakeLists.txt}
für ein ROS2-Paket tut.
\end{aufgabeBox}

\section{Kompiliert vs. interpretiert}

\begin{wissenBox}
Python übersetzt deinen Code Zeile für Zeile beim Ausführen (Interpreter).
C++ übersetzt deinen Code \emph{vor} der Ausführung in Maschinencode (Compiler).

\begin{center}
\begin{tabular}{lll}
\hline
\textbf{} & \textbf{Python} & \textbf{C++} \\
\hline
Ausführung & Interpreter zur Laufzeit & direkte CPU-Instruktionen \\
Fehler & zur Laufzeit sichtbar & schon beim Kompilieren \\
Geschwindigkeit & 10--100$\times$ langsamer & Referenzwert \\
Typfehler & zur Laufzeit (ohne mypy) & immer beim Kompilieren \\
Speicher & automatisch (GC) & manuell / Smart Pointer \\
\hline
\end{tabular}
\end{center}

Für \texttt{motor\_node}: Python-Overhead bei der Serial-Kommunikation
spielt kaum eine Rolle. Aber für einen Impedanzregler,
der 1000 Mal pro Sekunde Matrix-Operationen durchführt,
ist C++ der einzig richtige Weg.
\end{wissenBox}

\subsection{Drei Schritte vom Quelltext zum Programm}

\begin{lstlisting}[style=cpp]
// hallo.cpp -- das kleinste C++-Programm
#include <iostream>

int main() {
    std::cout << "Hallo ROS2-Welt!" << std::endl;
    return 0;
}
\end{lstlisting}

\begin{lstlisting}[language=bash]
# Schritt 1: Kompilieren (Quelltext -> Objektdatei)
g++ -c hallo.cpp -o hallo.o

# Schritt 2: Linken (Objektdatei + Bibliotheken -> Programm)
g++ hallo.o -o hallo

# Oder in einem Schritt:
g++ -std=c++17 hallo.cpp -o hallo
./hallo
# Hallo ROS2-Welt!
\end{lstlisting}

\subsection{CMake -- der Build-Orchestrator}

\begin{lstlisting}[style=cmake]
# CMakeLists.txt -- minimal
cmake_minimum_required(VERSION 3.8)
project(hallo_welt)

# C++17 einschalten
set(CMAKE_CXX_STANDARD 17)
set(CMAKE_CXX_STANDARD_REQUIRED ON)

# Executable definieren
add_executable(hallo hallo.cpp)
\end{lstlisting}

\begin{lstlisting}[language=bash]
# CMake-Build
mkdir build && cd build
cmake ..
make
./hallo
\end{lstlisting}

\subsection{ament\_cmake -- CMake für ROS2}

\texttt{ament\_cmake} ist ein dünner Wrapper über CMake,
der ROS2-Pakete als Abhängigkeiten kennt
und \texttt{colcon} als Build-Frontend nutzt.

\begin{lstlisting}[style=cmake]
# CMakeLists.txt für ein ROS2-C++-Paket
cmake_minimum_required(VERSION 3.8)
project(robotik_uno_q_cpp)

set(CMAKE_CXX_STANDARD 17)

# ament_cmake: Basis-ROS2-Build-System
find_package(ament_cmake REQUIRED)
# rclcpp: C++ ROS2-Bibliothek
find_package(rclcpp REQUIRED)
# Nachrichtentypen
find_package(geometry_msgs REQUIRED)
find_package(std_srvs REQUIRED)

# Executable
add_executable(motor_node src/motor_node.cpp)
ament_target_dependencies(motor_node
    rclcpp geometry_msgs std_srvs)

# Installationspfad
install(TARGETS motor_node
    DESTINATION lib/${PROJECT_NAME})

ament_package()
\end{lstlisting}

\begin{loesungBox}
\begin{lstlisting}[language=bash]
# Schritt-für-Schritt: erstes ROS2-C++-Paket
cd ~/ros2_ws/src
ros2 pkg create --build-type ament_cmake robotik_uno_q_cpp \
    --dependencies rclcpp geometry_msgs std_srvs

# Zeigt die erzeugte Struktur:
tree robotik_uno_q_cpp/
# robotik_uno_q_cpp/
# ├── CMakeLists.txt
# ├── package.xml
# ├── include/robotik_uno_q_cpp/
# └── src/

colcon build --packages-select robotik_uno_q_cpp
source install/setup.bash
\end{lstlisting}
\end{loesungBox}

\begin{projektBox}[robotik-uno-q -- Projektstand]
\textbf{Heute:} C++-Paket \texttt{robotik\_uno\_q\_cpp} erstellt und gebaut.\\
\textbf{Nächstes Kapitel:} Smart Pointer und Klassen -- Grundlage für rclcpp.
\end{projektBox}

\begin{merksatzBox}
\textbf{Merksatz:} \texttt{find\_package(rclcpp REQUIRED)} holt rclcpp.
\texttt{ament\_target\_dependencies()} verknüpft Abhängigkeiten mit dem Executable.
\texttt{ament\_package()} muss die letzte Zeile sein.
\texttt{colcon build} baut, \texttt{source install/setup.bash} macht den Node sichtbar.
\end{merksatzBox}

\begin{fehlerBox}
\begin{itemize}
\item \texttt{find\_package} vor \texttt{ament\_cmake}: Reihenfolge erzwungen --
  \texttt{ament\_cmake} muss als Erstes kommen.
\item \texttt{ament\_package()} fehlt: \texttt{colcon} findet das Paket nicht.
\item Header-Dateien in \texttt{src/} statt \texttt{include/}: funktioniert,
  aber verletzt ROS2-Konvention.
\end{itemize}
\end{fehlerBox}

\begin{haubeBox}
\textbf{Warum colcon statt make direkt?}

\texttt{colcon} (Collective Construction) verwaltet
Abhängigkeiten zwischen Paketen im Workspace.
Es baut Pakete in der richtigen Reihenfolge --
erst \texttt{rclcpp}, dann dein Paket, das rclcpp braucht.
\texttt{make} kennt nur das eine Paket, das du gerade baust.
In einem Workspace mit 20 Paketen wäre manuelles Reihenfolge-Management
ein Albtraum.
\end{haubeBox}

\begin{robotikBox}
Das offizielle \texttt{ros2\_control}-Paket
(Industriestandard für Motorsteuerung) ist in C++ geschrieben
und nutzt exakt dieses CMake-Muster.
Wer \texttt{CMakeLists.txt} für \texttt{motor\_node} versteht,
versteht auch, wie \texttt{ros2\_control} gebaut wird.
\end{robotikBox}


% ─────────────────────────────────────────────────────────────
\chapter{Wir wollen Speicher sicher verwalten –
  Smart Pointer und RAII}
% ─────────────────────────────────────────────────────────────

\begin{lernziele}
Du verstehst, warum rclcpp fast ausschließlich
\texttt{std::shared\_ptr} verwendet,
kannst RAII erklären und weißt,
warum \texttt{new}/\texttt{delete} in modernem ROS2-Code nicht auftaucht.
\end{lernziele}

\section{Situation}

rclcpp-Code sieht so aus:

\begin{lstlisting}[style=cpp]
auto node = std::make_shared<MotorNode>();
rclcpp::spin(node);
\end{lstlisting}

Warum \texttt{make\_shared}? Was ist ein \texttt{shared\_ptr}?
Und warum nicht einfach \texttt{new MotorNode()}?

\begin{aufgabeBox}
Verstehe Smart Pointer (\texttt{shared\_ptr}, \texttt{unique\_ptr})
und das RAII-Prinzip, damit du rclcpp-Code lesen und schreiben kannst.
\end{aufgabeBox}

\section{Das Problem mit rohen Zeigern}

\begin{wissenBox}
In klassischem C++ reservierst du Speicher mit \texttt{new}
und gibst ihn mit \texttt{delete} wieder frei.
Das ist fehleranfällig:

\begin{lstlisting}[style=cpp]
// Klassisch -- gefährlich
MotorNode* node = new MotorNode();
rclcpp::spin(node);
// Was, wenn spin() eine Exception wirft?
delete node;  // wird eventuell nie erreicht -- Memory Leak!
\end{lstlisting}

\textbf{Smart Pointer lösen das Problem:}

\begin{lstlisting}[style=cpp]
// Modern -- sicher
auto node = std::make_shared<MotorNode>();
rclcpp::spin(node);
// Wenn spin() endet oder eine Exception wirft:
// node's Destruktor läuft automatisch -- kein Memory Leak
\end{lstlisting}
\end{wissenBox}

\subsection{shared\_ptr -- gemeinsamer Besitz}

\begin{lstlisting}[style=cpp]
#include <memory>
#include <iostream>

class Serial {
public:
    Serial(const std::string& port) {
        std::cout << "Port " << port << " geoeffnet\n";
    }
    ~Serial() {
        std::cout << "Port geschlossen\n";
    }
};

{   // Neuer Scope
    auto ser1 = std::make_shared<Serial>("/dev/arduino");
    // Referenzzähler: 1
    {
        auto ser2 = ser1;   // Kopie des shared_ptr
        // Referenzzähler: 2
        // ser2 wird hier zerstört:
    }
    // Referenzzähler: 1 -- Serial lebt noch
}
// Referenzzähler: 0 -- ~Serial() wird aufgerufen
\end{lstlisting}

rclcpp verwendet \texttt{shared\_ptr} überall,
weil Nodes, Publisher und Subscriber
von mehreren Stellen gleichzeitig referenziert werden.

\subsection{unique\_ptr -- exklusiver Besitz}

\begin{lstlisting}[style=cpp]
// unique_ptr: nur ein Besitzer
auto serial = std::make_unique<Serial>("/dev/arduino");
// serial kann nicht kopiert werden, nur verschoben:
auto serial2 = std::move(serial);  // serial ist jetzt nullptr
\end{lstlisting}

\textbf{Faustregel:}
\texttt{unique\_ptr} für interne Ressourcen eines Nodes
(Serial-Port, Puffer).
\texttt{shared\_ptr} für Nodes selbst (rclcpp erwartet das).

\subsection{RAII -- Resource Acquisition Is Initialization}

\begin{lstlisting}[style=cpp]
class SerialPort {
public:
    // Konstruktor: Ressource öffnen
    explicit SerialPort(const std::string& port, int baud)
        : port_(port) {
        // Hier: Serial öffnen
        offen_ = true;
        std::cout << "Port " << port << " @ " << baud << " Baud\n";
    }
    // Destruktor: Ressource schließen -- automatisch
    ~SerialPort() {
        if (offen_) {
            std::cout << "Port " << port_ << " geschlossen\n";
        }
    }
    // Kein Kopieren erlaubt
    SerialPort(const SerialPort&) = delete;
    SerialPort& operator=(const SerialPort&) = delete;

private:
    std::string port_;
    bool offen_{false};
};
\end{lstlisting}

RAII bedeutet: Die Ressource wird im Konstruktor geöffnet
und im Destruktor geschlossen -- automatisch, auch bei Exceptions.

\begin{projektBox}[robotik-uno-q -- Projektstand]
\textbf{Heute:} RAII-Konzept und Smart Pointer verstanden --
Grundlage für den \texttt{MotorNode}.\\
\textbf{Nächstes Kapitel:} Lambda, Templates und modernes C++17.
\end{projektBox}

\begin{merksatzBox}
\textbf{Merksatz:}
Kein \texttt{new}/\texttt{delete} in modernem C++.
\texttt{make\_shared<T>(args)} für geteilten Besitz (rclcpp-Nodes).
\texttt{make\_unique<T>(args)} für exklusiven Besitz (interne Ressourcen).
RAII: Konstruktor öffnet, Destruktor schließt -- immer, auch bei Exceptions.
\end{merksatzBox}

\begin{fehlerBox}
\begin{itemize}
\item \texttt{shared\_ptr} auf Stack-Objekte: Destruktor läuft zweimal.
  Immer \texttt{make\_shared} verwenden.
\item Zirkuläre \texttt{shared\_ptr}: A hält B, B hält A → Memory Leak.
  Lösung: schwache Referenz (\texttt{weak\_ptr}) für eine Richtung.
\item \texttt{unique\_ptr} kopieren: Compiler-Fehler.
  \texttt{std::move(ptr)} für Eigentumsübertragung.
\end{itemize}
\end{fehlerBox}

\begin{haubeBox}
\textbf{Warum rclcpp auf shared\_ptr setzt}

rclcpp muss Nodes, Publisher und Subscriber registrieren
und intern referenzieren -- auch nachdem der Aufrufer
seinen Pointer verloren hat.
\texttt{shared\_ptr} garantiert: Der Node lebt, solange rclcpp ihn braucht.
Python löst das mit dem Garbage Collector.
C++ löst es mit \texttt{shared\_ptr} -- explizit, aber vorhersagbar.
\end{haubeBox}


% ─────────────────────────────────────────────────────────────
\chapter{Wir wollen modernen C++-Code lesen –
  auto, Lambda und Templates}
% ─────────────────────────────────────────────────────────────

\begin{lernziele}
Du kannst \texttt{auto}, Lambda-Ausdrücke und Template-Funktionen
in rclcpp-Code lesen und schreiben --
diese drei Features tauchen in fast jeder rclcpp-Zeile auf.
\end{lernziele}

\section{Situation}

rclcpp-Code enthält oft Zeilen wie:

\begin{lstlisting}[style=cpp]
auto sub = this->create_subscription<sensor_msgs::msg::Range>(
    "/sensor/rechts", 10,
    [this](const sensor_msgs::msg::Range::SharedPtr msg) {
        this->_cb_rechts(msg);
    });
\end{lstlisting}

Drei Dinge auf einmal: \texttt{auto}, Template mit \texttt{<>}
und Lambda mit \texttt{[this](...)\{...\}}.
Ohne Verständnis dieser Features ist rclcpp-Code unlesbar.

\begin{aufgabeBox}
Verstehe \texttt{auto}, Lambda und Templates so weit,
dass du jeden rclcpp-Subscriber und -Publisher selbst schreiben kannst.
\end{aufgabeBox}

\section{auto -- Typ-Inferenz}

\begin{lstlisting}[style=cpp]
// Ohne auto: ausführlich
std::shared_ptr<rclcpp::Publisher<geometry_msgs::msg::Twist>> pub;

// Mit auto: kompakter, aber identisch
auto pub = this->create_publisher<geometry_msgs::msg::Twist>(
    "/cmd_vel", 10);
// Der Compiler leitet den Typ aus dem Rückgabewert ab
\end{lstlisting}

\texttt{auto} ist kein dynamischer Typ wie Python's Typsystem.
Der Typ wird beim Kompilieren fest fixiert --
\texttt{auto} ist syntaktischer Zucker,
kein Laufzeit-Feature.

\section{Lambda -- anonyme Funktionen}

\begin{lstlisting}[style=cpp]
// Lambda-Syntax:  [capture](parameter) -> rückgabetyp { body }
auto add = [](int a, int b) -> int { return a + b; };
add(3, 4);   // -> 7

// [this]: Zugriff auf Klassenmitglieder
class MotorNode : public rclcpp::Node {
    void _richte_subscriber_ein() {
        sub_ = this->create_subscription<geometry_msgs::msg::Twist>(
            "/cmd_vel", 10,
            [this](const geometry_msgs::msg::Twist::SharedPtr msg) {
                // 'this' ist verfügbar dank [this]-Capture
                this->_verarbeite_twist(msg);
            });
    }
};
\end{lstlisting}

\section{Templates -- generischer Code}

\begin{lstlisting}[style=cpp]
// Template-Funktion: funktioniert für jeden Typ T
template<typename T>
T klemme(T wert, T min_val, T max_val) {
    return std::max(min_val, std::min(wert, max_val));
}

klemme(87, 0, 100);   // int: 87
klemme(4.5f, 0.0f, 4.0f);  // float: 4.0

// rclcpp-Publisher ist eine Template-Klasse:
// create_publisher<T>(topic, qos)
// T ist der Nachrichtentyp -- wird zur Kompilierzeit eingesetzt
auto pub = create_publisher<geometry_msgs::msg::Twist>("/cmd_vel", 10);
\end{lstlisting}

\begin{projektBox}[robotik-uno-q -- Projektstand]
\textbf{Heute:} \texttt{auto}, Lambda und Templates lesbar --
rclcpp-Code wird verständlich.\\
\textbf{Nächstes Kapitel:} Multithreading -- \texttt{std::thread} und Mutexe.
\end{projektBox}

\begin{merksatzBox}
\textbf{Merksatz:}
\texttt{auto} spart Tipparbeit, ändert aber nichts am Typ.
Lambda \texttt{[this](...)\{...\}} ist eine anonyme Funktion
mit Zugriff auf \texttt{this}.
\texttt{template<typename T>} macht Code für beliebige Typen verwendbar.
Alle drei sind in rclcpp allgegenwärtig.
\end{merksatzBox}

\begin{fehlerBox}
\begin{itemize}
\item \texttt{[this]} vergessen: Zugriff auf Klassenmitglieder
  ergibt Compiler-Fehler (\texttt{this} nicht im Scope des Lambda).
\item \texttt{auto} bei Referenzen:
  \texttt{auto x = obj.get\_wert()} kopiert den Wert.
  \texttt{auto\& x = ...} ist die Referenz.
\item Template-Fehler: Compiler-Meldungen bei Templates sind
  oft sehr lang -- immer nur die erste Zeile lesen,
  die den eigentlichen Fehler nennt.
\end{itemize}
\end{fehlerBox}


% ─────────────────────────────────────────────────────────────
\chapter{Wir wollen parallel arbeiten –
  Multithreading in rclcpp}
% ─────────────────────────────────────────────────────────────

\begin{lernziele}
Du verstehst \texttt{std::thread}, \texttt{std::mutex} und
\texttt{std::atomic}, weißt wann der \texttt{MultiThreadedExecutor}
nötig ist und implementierst den Serial-Lese-Thread
im \texttt{motor\_node} korrekt.
\end{lernziele}

\section{Situation}

\texttt{motor\_node} soll gleichzeitig:
\begin{enumerate}
\item \texttt{/cmd\_vel}-Nachrichten empfangen (Subscriber-Callback)
\item Bestätigungen vom Arduino lesen (Serial-Polling)
\end{enumerate}

Mit \texttt{SingleThreadedExecutor} blockiert Serial-Polling
alle anderen Callbacks. Das braucht einen zweiten Thread.

\begin{aufgabeBox}
Implementiere einen Hintergrund-Thread für Serial-Polling
im \texttt{motor\_node}, der sicher auf gemeinsame Daten zugreift.
\end{aufgabeBox}

\section{std::thread und std::mutex}

\begin{lstlisting}[style=cpp]
#include <thread>
#include <mutex>
#include <atomic>
#include <string>

class MotorNode : public rclcpp::Node {
public:
    MotorNode() : Node("motor") {
        laufe_.store(true);
        lese_thread_ = std::thread(&MotorNode::_lese_arduino, this);
    }

    ~MotorNode() {
        laufe_.store(false);
        if (lese_thread_.joinable()) {
            lese_thread_.join();   // Warten bis Thread endet
        }
    }

private:
    std::atomic<bool> laufe_{false};  // thread-sicher ohne Mutex
    std::thread lese_thread_;
    std::mutex ser_mutex_;             // Schutz für Serial-Port
    std::string letzte_antwort_;

    void _lese_arduino() {
        while (laufe_.load()) {
            // Serial lesen (blockiert bis Zeile kommt)
            std::string zeile = _lese_zeile();
            {
                std::lock_guard<std::mutex> lock(ser_mutex_);
                letzte_antwort_ = zeile;
            }  // lock wird hier automatisch freigegeben (RAII)
        }
    }

    // Subscriber-Callback -- läuft im Executor-Thread
    void _cb_cmd_vel(const geometry_msgs::msg::Twist::SharedPtr msg) {
        // serial schreiben ist auch mutex-geschützt
        std::lock_guard<std::mutex> lock(ser_mutex_);
        _sende_befehl(msg->linear.x, msg->angular.z);
    }
};
\end{lstlisting}

\begin{wissenBox}
\textbf{std::atomic:} Für einfache Typen (\texttt{bool}, \texttt{int})
ist \texttt{atomic} schneller als \texttt{mutex} --
Lesen und Schreiben sind atomar (unteilbar) garantiert.

\textbf{std::lock\_guard:} RAII für Mutex --
\texttt{lock\_guard} sperrt den Mutex im Konstruktor,
gibt ihn im Destruktor frei -- auch bei Exceptions.
Kein \texttt{mutex.unlock()} nötig.

\textbf{Deadlock:} Wenn Thread~A Mutex~1 hält und Mutex~2 wartet,
und Thread~B Mutex~2 hält und Mutex~1 wartet --
beide warten ewig. Lösung: Mutexe immer in der gleichen Reihenfolge sperren.
\end{wissenBox}

\begin{projektBox}[robotik-uno-q -- Projektstand]
\textbf{Heute:} Multithreading-Grundlagen verstanden.\\
\textbf{Teil~II:} rclcpp -- die eigentlichen ROS2-Nodes.
\end{projektBox}

\begin{merksatzBox}
\textbf{Merksatz:}
\texttt{std::atomic<bool>} für einfache Flags zwischen Threads.
\texttt{std::lock\_guard<std::mutex>} für alle anderen gemeinsamen Daten.
Threads immer mit \texttt{join()} sauber beenden.
\texttt{MultiThreadedExecutor} für rclcpp-Nodes mit parallelen Callbacks.
\end{merksatzBox}

\begin{fehlerBox}
\begin{itemize}
\item Shared Data ohne Mutex: Data Race -- undefiniertes Verhalten,
  nicht reproduzierbar, schwerste Kategorie von Bugs.
\item Thread-Destruktor ohne \texttt{join()}: \texttt{std::terminate()} --
  das Programm bricht ab.
\item \texttt{mutex.lock()} ohne \texttt{unlock()} (statt \texttt{lock\_guard}):
  Deadlock beim zweiten Aufruf.
\end{itemize}
\end{fehlerBox}
